\chapter{Strings on K3: Moduli Space and Dualities}\label{ch:stringsk3}

Two strands of this book meet in this chapter. In Parts~II and~III we studied strings on
tori: the momentum and winding numbers form an even self-dual lattice $\Gamma^{d,d}$, the
background fields $G$ and $B$ select a positive-definite $d$-plane in $\R^{d,d}$, and the
duality group $\Odd{d}$ is the isometry group of the lattice (\cref{ch:narain,ch:odd}). In
Part~IV we studied a purely geometric object, the K3 surface, and found that its cohomology
carries the even unimodular lattice $\Gamma^{3,19}=3U\oplus2E_8(-1)$ and, after adding
$H^0\oplus H^4$, the Mukai lattice $\Gamma^{4,20}$ (\cref{ch:k3lattice,ch:mukai}). The
question of this chapter is: \emph{what happens when a string propagates on a K3 surface?}

The answer has the same shape as for a torus. The string sees, besides the Ricci-flat
metric, a $B$-field; together they are $80$ real parameters, and they assemble into the
choice of a positive-definite \emph{four}-plane $\Pi\subset\R^{4,20}$. The moduli space is
\begin{equation}\label{eq:stringsk3:main}
  \mathcal M_\sigma \;=\; O(\Gamma^{4,20})\,\backslash\, O(4,20)\,/\,\big(O(4)\times O(20)\big),
\end{equation}
which is literally a Narain moduli space --- the one of the \emph{heterotic} string on a
four-torus. This coincidence is the first evidence for a non-perturbative duality between
two very different string theories, type~IIA on K3 and heterotic on $T^4$, and it predicts
that a type~IIA string acquires non-abelian gauge symmetry when a two-sphere inside the K3
surface shrinks to zero size. A student should care for two reasons: this is the simplest
example in which a curved compactification is controlled exactly by a lattice, and it is
the six-dimensional parent of the $\KT$ compactification that the rest of the book is about
(\cref{ch:k3t2}).

Almost everything in this chapter is Tier~\TL: it is taken from Aspinwall's lectures, with
the derivations filled in. The machine-checked content is small and precisely delimited:
signatures, ranks and parity of the lattices involved (\cref{sec:stringsk3:lean}). Our
source makes its own assumptions explicit, and we repeat them where they enter.

\section{The sigma model with K3 target}\label{sec:stringsk3:sigma}

\subsection{Action and conformal invariance}
Let $X$ be the target manifold and $x:\Sigma_{\rm ws}\to X$ the embedding of the worldsheet
$\Sigma_{\rm ws}$ (a Riemann surface; we write $\Sigma_{\rm ws}$ because $\Sigma$ will
denote a three-plane below). In conformal gauge the bosonic part of the action is, in
Aspinwall's normalization,
\begin{equation}\label{eq:stringsk3:action}
  S=\frac{\ii}{8\pi\ap}\int_{\Sigma_{\rm ws}}\big(g_{ij}-B_{ij}\big)\,
     \partial x^i\,\bar\partial x^j\,\dd^2z\;-\;\frac{1}{2\pi}\int_{\Sigma_{\rm ws}}\Phi\,R^{(2)}\,\dd^2z+\dots
\end{equation}
\source{Asp §3.1, eq.~(60), ll.~1559--1576}. Here $g_{ij}$ is a Riemannian metric on $X$,
$B_{ij}$ the components of a real two-form $B$, $\Phi$ the dilaton and $R^{(2)}$ the scalar
curvature of the worldsheet; the dots are fermionic terms. This is the same sigma model as
in \cref{ch:backgrounds}, written in complex worldsheet coordinates; the overall numerical
factors differ from Tong's normalization and play no role below, so we do not convert them
(the coefficient of the dilaton term is garbled in our copy of the source; only its
structure, $\Phi$ times the Euler density, is used).

For the two-dimensional theory to define a string background it must be conformally
invariant. As in \cref{ch:backgrounds}, this is a condition order by order in $\ap/R^2$, where
$R$ is a typical radius of $X$. The simplest solution at leading order is: $\Phi$ constant,
$\dd B=0$, and $g_{ij}$ \emph{Ricci-flat} \source{Asp §3.1, ll.~1583--1592}. The two
compact four-dimensional examples used in this book are the flat torus $T^4$ and the K3
surface with its Yau metric (\cref{ch:k3}).

\subsection{Worldsheet supersymmetry and hyperk\"ahler geometry}
Add the worldsheet fermions $\psi^i$, which transform as sections of the (pulled-back)
cotangent bundle of $X$. A supersymmetry variation has the form
$\delta_\epsilon x^i=\bar\epsilon\,l^i{}_j\psi^j$. With one supersymmetry, $l$ can be set to
the identity. A second supersymmetry requires a second tensor $l$, which turns out to be an
almost complex structure making $X$ K\"ahler; $N=4$ requires three complex structures, that is,
a hyperk\"ahler manifold, and conversely \source{Asp §3.1, ll.~1601--1611}. A K3 surface with
a Ricci-flat metric has holonomy $SU(2)\cong Sp(1)$ and carries a two-sphere of complex
structures $aI+bJ+cK$, $a^2+b^2+c^2=1$ (\cref{ch:k3}). Because a conformal field theory splits
into left- and right-movers, each sector carries its own $N=4$ algebra: \emph{the sigma model
on a K3 surface is an $N=(4,4)$ superconformal field theory.}
\index{N=(4,4) superconformal field theory@$N=(4,4)$ superconformal field theory}
This has a consequence that distinguishes K3 from Calabi--Yau threefolds: for $N=4$ sigma
models there are no perturbative corrections to the Ricci-flat metric beyond leading order,
and there are arguments that this persists non-perturbatively, so the Ricci-flat metric is an
exact solution \source{Asp §3.1, ll.~1618--1625} \TL. The central charge is $c=6$ on each side
(four bosons and four fermions, $4+4\cdot\tfrac12$; \cref{ch:cft}), and the $N=4$ algebra
contains an affine $\mathfrak{su}(2)$. We will meet its characters again in
\cref{ch:ellgenus}.

\begin{physicsbox}[Physical picture: why the torus analogy is not an accident]
For a torus the metric is flat and the zero modes of the string are labelled by a lattice.
For K3 the metric is curved and unknown in closed form,
yet $N=(4,4)$ supersymmetry is strong enough to protect the theory from corrections. The
price is that we cannot solve the model at a generic point; what we \emph{can} determine is
the space of all such models. The logic is that of \cref{ch:k3lattice}: there, Yau's theorem
allowed us to replace ``a Ricci-flat metric'' by ``a positive three-plane in $\R^{3,19}$''
without ever writing the metric down. Here the same trick is played one dimension higher.
\end{physicsbox}

\section{Counting moduli: \texorpdfstring{$58+22=80$}{58+22=80}}\label{sec:stringsk3:count}

\subsection{Three sets of parameters}
The action \eqref{eq:stringsk3:action} contains three kinds of data. We ask, for each, how
many deformations change the conformal field theory.

\paragraph{The metric.} It must be Ricci-flat. The moduli space of Einstein metrics on a K3
surface (orbifold points included) is
\begin{equation}\label{eq:stringsk3:einstein}
  \mathcal M_E\cong O(\Gamma^{3,19})\backslash O(3,19)/\big(O(3)\times O(19)\big)\times\R_+ ,
\end{equation}
the Grassmannian of positive three-planes $\Sigma\subset\R^{3,19}=H^2(X,\R)$, modulo lattice
isometries, times the volume \source{Asp §2.4, Theorem~3, eqs.~(39)--(40), ll.~875--893}
(\cref{ch:k3lattice}). The plane $\Sigma$ is spanned by the real and imaginary parts of the
holomorphic two-form and by the K\"ahler form $J$. The dimension of $\mathcal M_E$ is
$3\times19+1=58$.
Assuming every inequivalent Ricci-flat metric gives an inequivalent conformal field theory,
the metric contributes $58$ parameters \source{Asp §3.2, ll.~1639--1642}.

\paragraph{The $B$-field.} It enters through $\int_{x(\Sigma_{\rm ws})}B$, the integral of $B$
over the image of the worldsheet, a closed two-cycle. By Stokes' theorem an exact piece
$B=\dd\Lambda$ integrates to zero, so only the cohomology class $[B]\in H^2(X,\R)$ matters.
Since $b_2(K3)=22$, the $B$-field contributes $22$ parameters
\source{Asp §3.2, ll.~1643--1654}.\index{B-field@$B$-field!on K3}

\paragraph{The dilaton.} For constant $\Phi$ the last term of \eqref{eq:stringsk3:action} is
topological: by Gauss--Bonnet it contributes $2\Phi(g-1)=-\Phi\,\chi(\Sigma_{\rm ws})$ on a
worldsheet of genus $g$. A genus-$g$ amplitude is therefore weighted by $\lambda^{2g-2}$ with
the string coupling $\lambda=\ee^{\Phi}$, and nothing changes in the conformal field theory on
a fixed worldsheet \source{Asp §3.2, ll.~1655--1667}. It is
\emph{not} a modulus of the conformal field theory. It returns in
\cref{sec:stringsk3:duality} as a modulus of the string theory.

\begin{equation}\label{eq:stringsk3:80}
  \dim_\R\mathcal M_\sigma \;=\; 58+22 \;=\; 80 \;=\; 4\times20 .
\end{equation}
\source{Asp §3.2, l.~1670}. The factorization $80=4\cdot20$ is the first hint of the answer.

\subsection{The local form of the moduli space}
To find the space itself, Aspinwall uses a holonomy argument, which we retrace.

\emph{Step 1.} The symmetric space $O(3,19)/(O(3)\times O(19))$ in
\eqref{eq:stringsk3:einstein} has holonomy algebra $\mathfrak{so}(3)\oplus\mathfrak{so}(19)$. The
$\mathfrak{so}(3)\cong\mathfrak{su}(2)$ rotates the three-plane $\Sigma$ into itself, i.e.\ it
rotates the sphere of complex structures.

\emph{Step 2.} In a conformal field theory left- and right-movers are independent, so each
side has its own $\mathfrak{su}(2)$; indeed each $N=4$ algebra contains an affine
$\mathfrak{su}(2)$. Therefore the $\mathfrak{so}(3)$ is promoted to
$\mathfrak{su}(2)\oplus\mathfrak{su}(2)\cong\mathfrak{so}(4)$, which acts on the marginal
operators, the tangent vectors of the moduli space
\source{Asp §3.2, ll.~1674--1690}.

\emph{Step 3.} We need an $80$-dimensional Riemannian manifold whose holonomy contains
$SO(4)$ acting non-trivially on every tangent direction, and which is non-compact (a K3
surface can be made arbitrarily large). The classification of holonomy groups by Berger and
Simons then leaves one possibility \source{Asp §3.2, Theorem~5, ll.~1691--1713}:

\begin{theorem}[Local form; Aspinwall's Theorem~5, \TL]\label{thm:stringsk3:local}
Assuming that the deformations of $g$ and $B$ counted above are effective, every smooth
neighbourhood in the moduli space of $N=(4,4)$ sigma models with K3 target is isomorphic to an
open subset of
\begin{equation}\label{eq:stringsk3:teich}
  \mathcal T_\sigma=\frac{O(4,20)}{O(4)\times O(20)} .
\end{equation}
\end{theorem}

The dimension can be checked directly. The group $O(p,q)$ has the same dimension as
$O(p+q)$, namely $\tfrac12(p+q)(p+q-1)$, so
\begin{equation}\label{eq:stringsk3:dimpq}
  \dim\frac{O(p,q)}{O(p)\times O(q)}
  =\frac{(p+q)(p+q-1)}2-\frac{p(p-1)}2-\frac{q(q-1)}2=pq ,
\end{equation}
and for $(p,q)=(4,20)$ this is $276-6-190=80$. A more geometric way to see $pq$: the coset is
the Grassmannian of positive-definite $p$-planes $\Pi\subset\R^{p,q}$; a nearby plane is the
graph of a linear map $\Pi\to\Pi^\perp$, and such maps form a vector space of dimension $pq$.
This is exactly how $G+B$ appeared as a map in \cref{ch:narain}.\index{Grassmannian}

\begin{pitfallbox}[Common pitfall: $80$ is not ``$58$ plus $22$ Ramond--Ramond fields'']
The $22$ extra moduli are periods of the Neveu--Schwarz $B$-field, present in every closed
string theory. They are not Ramond--Ramond fields: type~IIA has R--R potentials of odd degree
only (a one-form and a three-form), a K3 surface has no odd cycles, and so the R--R sector
contributes \emph{no} scalars on K3 \source{Asp §4.2, ll.~2348--2353}. Nor is the dilaton among
the $80$: it is an $81$st modulus, see \eqref{eq:stringsk3:MIIA}. We stress this because the
informal comment attached to the Lean declaration \lean{witten\_moduli\_dim\_is\_80}
(\cref{sec:stringsk3:lean}) gets precisely this point wrong; the theorem itself, $4\times20=80$,
is unaffected.
\end{pitfallbox}

\section{Geometry of the four-plane}\label{sec:stringsk3:fourplane}

A point of $\mathcal T_\sigma$ is a positive-definite four-plane $\Pi\subset\R^{4,20}$. How does
one read off a metric, a $B$-field and a volume from $\Pi$? The answer requires a choice, and
the choice is the origin of the stringy dualities.

\subsection{Choosing a null vector}
Let $\Gamma^{4,20}\subset\R^{4,20}$ be the even unimodular lattice of signature $(4,20)$; by the
classification of \cref{ch:lattices} it is unique, $\Gamma^{4,20}\cong4U\oplus2E_8(-1)
\cong\Gamma^{3,19}\oplus U$. Fix a primitive vector $w\in\Gamma^{4,20}$ with $w\cdot w=0$, and a
second null vector $w^*$ with $w\cdot w^*=1$. Then $w$ and $w^*$ span a copy of $U$ and
\begin{equation}\label{eq:stringsk3:wperp}
  \frac{w^\perp}{w}\cap\Gamma^{4,20}\;\cong\;\Gamma^{3,19},\qquad
  \R^{4,20}=\R w\oplus\R w^*\oplus\R^{3,19}
\end{equation}
\source{Asp §3.3, eq.~(65), ll.~1741--1753}. Here $w^\perp$ is the $23$-dimensional space of
vectors orthogonal to $w$; it contains $w$ itself because $w$ is null, and dividing by $w$
leaves $22$ dimensions. In the language of \cref{ch:mukai}, $\R^{3,19}=H^2(X,\R)$, $w$ generates
$H^4(X,\Z)$ and $w^*$ generates $H^0(X,\Z)$ (up to the sign convention for the Mukai pairing
discussed there).

\subsection{From \texorpdfstring{$\Pi$}{Pi} to \texorpdfstring{$(\Sigma,B,V)$}{(Sigma,B,V)}
and back}
Given $\Pi$, set $\Sigma'=\Pi\cap w^\perp$. Since $\Pi$ is four-dimensional and $w^\perp$ has
codimension one, $\Sigma'$ is a three-plane (it cannot be all of $\Pi$: otherwise $w$ would be
orthogonal to a maximal positive-definite subspace, hence lie in the negative-definite
$\Pi^\perp$, contradicting $w^2=0$, $w\neq0$). Project $\Sigma'$ to $w^\perp/w=\R^{3,19}$ to
obtain a positive three-plane $\Sigma$: this is the Einstein metric of fixed volume. Let
$B'\in\Pi$ be the vector orthogonal to $\Sigma'$ with $B'\cdot w=1$, and decompose it as
\begin{equation}\label{eq:stringsk3:Bprime}
  B'=\alpha\,w+w^*+B,\qquad \alpha\in\R,\quad B\in\R^{3,19}=H^2(X,\R)
\end{equation}
\source{Asp §3.3, eq.~(66), ll.~1755--1768}. The component $B$ is the $B$-field. Let us make
the inverse map explicit, which the source leaves to the reader.

\begin{proposition}[The four-plane in coordinates]\label{thm:stringsk3:plane}
Let $\sigma_1,\sigma_2,\sigma_3$ be an orthonormal basis of a positive three-plane
$\Sigma\subset\R^{3,19}$, let $B\in\R^{3,19}$ and $\alpha\in\R$ with $2\alpha+B\cdot B>0$. Then
\begin{equation}\label{eq:stringsk3:basis}
  \xi_i=\sigma_i-(\sigma_i\cdot B)\,w\quad(i=1,2,3),\qquad \xi_4=B'=\alpha w+w^*+B
\end{equation}
are mutually orthogonal, with $\xi_i\cdot\xi_j=\delta_{ij}$ and
\begin{equation}\label{eq:stringsk3:volume}
  B'\cdot B'=2\alpha+B\cdot B .
\end{equation}
They span a positive-definite four-plane $\Pi$ with $\Pi\cap w^\perp=\mathrm{span}(\xi_1,
\xi_2,\xi_3)$.
\end{proposition}
\begin{proof}
Use $w^2=(w^*)^2=0$, $w\cdot w^*=1$ and $w\cdot\sigma=w^*\cdot\sigma=0$ for
$\sigma\in\R^{3,19}$. Then $\xi_i\cdot\xi_j=\sigma_i\cdot\sigma_j=\delta_{ij}$, because every
term containing $w$ pairs to zero. Next
$\xi_i\cdot B'=\sigma_i\cdot B-(\sigma_i\cdot B)\,(w\cdot w^*)=0$: the correction
$-(\sigma_i\cdot B)w$ is exactly what is needed to make $\xi_i$ orthogonal to $B'$. Finally
$B'\cdot B'=2\alpha\,(w\cdot w^*)+B\cdot B$. Since $\xi_i\cdot w=0$ and $B'\cdot w=1$, the
intersection with $w^\perp$ is spanned by the $\xi_i$.
\end{proof}

Aspinwall identifies the norm \eqref{eq:stringsk3:volume} with the volume of the K3 surface,
\begin{equation}\label{eq:stringsk3:vol}
  V=\mathrm{Vol}(X)=\int_XJ\wedge J=2\alpha+B\cdot B
\end{equation}
(in units of $\ap{}^2$), by matching the natural metrics on both sides of
\eqref{eq:stringsk3:decomp} below \source{Asp §3.3, eq.~(68), ll.~1786--1797 and footnote~9}
\TL. \Cref{thm:stringsk3:plane} makes this plausible: the condition that $\Pi$ be
positive-definite is precisely $V>0$. Counting parameters, $\Sigma$ has $3\times19=57$, $B$
has $22$ and $V$ has one:
\begin{equation}\label{eq:stringsk3:decomp}
  \frac{O(4,20)}{O(4)\times O(20)}\;\cong\;\frac{O(3,19)}{O(3)\times O(19)}\times\R^{22}\times\R_+ ,
  \qquad 80=57+22+1
\end{equation}
\source{Asp §3.3, eq.~(67), ll.~1770--1784}.\index{moduli space!of K3 sigma models}

\subsection{The discrete identifications}\label{sec:stringsk3:discrete}
Which isometries $g$ of the lattice are dualities, in the sense that $\Pi$ and $g\Pi$ define
the same theory? Those visible in the action \eqref{eq:stringsk3:action} fix $w$:
\begin{itemize}[leftmargin=*]
\item Diffeomorphisms of $X$ act on $H^2(X,\Z)$ by $O^+(\Gamma^{3,19})$ (\cref{ch:k3lattice});
  worldsheet complex conjugation, which reverses the sign of $B$, supplies the remaining
  element of $O(\Gamma^{3,19})$ \source{Asp §3.3, ll.~1802--1804 and 1823--1825}.
\item Integral shifts $B\mapsto B+e$, $e\in H^2(X,\Z)$: the action changes by
  $2\pi\ii\times$integer, so the path integral is unchanged
  \source{Asp §3.3, ll.~1804--1810}.\index{B-field@$B$-field!integral shift}
\end{itemize}
It is instructive to see that the $B$-shift really is an isometry of $\Gamma^{4,20}$.

\begin{proposition}[The $B$-shift as a lattice isometry]\label{thm:stringsk3:bshift}
For $e\in\Gamma^{3,19}$ define a linear map $T_e$ of $\R^{4,20}$ by
\begin{equation}\label{eq:stringsk3:Te}
  T_e(w)=w,\qquad T_e(x)=x-(x\cdot e)\,w\ \ (x\in\R^{3,19}),\qquad
  T_e(w^*)=w^*+e-\tfrac12(e\cdot e)\,w .
\end{equation}
Then $T_e$ preserves the inner product, maps $\Gamma^{4,20}$ to itself, and sends the plane
with data $(\Sigma,B,\alpha)$ to the plane with data
$\big(\Sigma,\;B+e,\;\alpha-B\cdot e-\tfrac12e\cdot e\big)$. In particular the volume
\eqref{eq:stringsk3:vol} is unchanged.
\end{proposition}
\begin{proof}
Isometry: $T_e(x)\cdot T_e(y)=x\cdot y$ since $w$ is null and orthogonal to $\R^{3,19}$;
$T_e(w^*)\cdot T_e(w)=1$; $T_e(w^*)\cdot T_e(x)=e\cdot x-(x\cdot e)=0$; and
$T_e(w^*)^2=e\cdot e-2\cdot\tfrac12e\cdot e=0$. Integrality: the only non-obvious coefficient
is $\tfrac12e\cdot e$, which is an integer \emph{because $\Gamma^{3,19}$ is even}. Action on
the plane: $T_e(B')=\alpha w+w^*+e-\tfrac12(e\cdot e)w+B-(B\cdot e)w$, which is of the form
\eqref{eq:stringsk3:Bprime} with the stated $B$ and $\alpha$; and
$T_e(\xi_i)=\sigma_i-(\sigma_i\cdot(B+e))w$. Finally
$2\alpha'+(B+e)^2=2\alpha-2B\cdot e-e^2+B^2+2B\cdot e+e^2=2\alpha+B^2$.
\end{proof}

The evenness of the lattice is therefore not a technicality: it is what makes the integral
$B$-shift a symmetry, exactly as for the torus, where $B_{ij}\mapsto B_{ij}+N_{ij}$ with $N$
integer antisymmetric lies in $\Odd{d}$ (\cref{ch:odd}). The maps $T_e$ together with
$O(\Gamma^{3,19})$ form the ``space group'' $O(\Gamma^{3,19})\ltimes\Gamma^{3,19}$
\source{Asp §3.3, eq.~(69), ll.~1826--1830}. These all fix $w$. One more element, moving $w$, is
supplied by mirror symmetry: for an algebraic K3 surface with transcendental lattice $\Lambda$
and ``quantum Picard lattice'' $\Upsilon=\mathrm{Pic}(X)\oplus U=\Lambda^\perp$, split $\Pi$ into
the two-planes $\Omega=\Pi\cap(\Lambda\otimes\R)$ (complex structure) and
$\mho=\Pi\cap(\Upsilon\otimes\R)$ (K\"ahler form and $B$-field); the mirror map exchanges
$(\Lambda,\Omega)\leftrightarrow(\Upsilon,\mho)$; in a Greene--Plesser example with $\Lambda\cong\Upsilon\cong\Gamma^{2,10}$ it is
realized by an element of $O(\Gamma^{4,20})$ exchanging two orthogonal $\Gamma^{2,10}$
sublattices, and this element together with the space group generates all of
$O(\Gamma^{4,20})$ \source{Asp §3.4, eqs.~(70)--(75), ll.~1860--1943}.
\index{mirror symmetry!for K3 surfaces}

\begin{theorem}[Global form; Aspinwall's Theorem~6, \TL]\label{thm:stringsk3:global}
Given \cref{thm:stringsk3:local}, completeness, and the assumption that the moduli space is
Hausdorff, the moduli space of conformally invariant sigma models on a K3 surface is
$\mathcal M_\sigma=O(\Gamma^{4,20})\backslash O(4,20)/(O(4)\times O(20))$
\source{Asp §3.4, Theorem~6, eq.~(76), ll.~1944--1950}.
\end{theorem}

The lattice $\Gamma^{4,20}$ is the total integral cohomology $H^*(X,\Z)$ with the pairing that
is non-zero only between degrees adding up to four --- the Mukai lattice of \cref{ch:mukai}
\source{Asp §3.4, ll.~1951--1957}. A point of $\mathcal M_\sigma$ fixes the conformal field
theory; only after choosing which null direction is called $H^0$ (equivalently $w=H^4$) can one
speak of ``the'' K3 geometry, and different choices can give very different-looking K3
surfaces with the same conformal field theory. Aspinwall calls this ``a kind of T-duality''
\source{Asp §3.4, ll.~1958--1963}.\index{T-duality!for K3 sigma models}

\section{The same coset from a torus: Narain revisited}\label{sec:stringsk3:narain}

The space \eqref{eq:stringsk3:main} has the form of a Narain moduli space. To exploit this we
need the torus story in the ``plane in a lattice'' language, allowing for different numbers
$n_L$, $n_R$ of left- and right-moving compact bosons.

\subsection{Planes and momenta}
Let $\Gamma^{n_L,n_R}\subset\R^{n_L,n_R}$ be an even self-dual lattice (which forces
$n_L-n_R\in8\Z$) and $\Pi$ a positive-definite $n_L$-plane. Every vector decomposes as
$p=p_L+p_R$ with $p_L\in\Pi$, $p_R\in\Pi^\perp$, so $p_L\cdot p_L\ge0$ and $p_R\cdot p_R\le0$.
The lattice points are the winding and momentum states, with left and right conformal weights
\begin{equation}\label{eq:stringsk3:weights}
  h=\tfrac12\,p_L\cdot p_L,\qquad \bar h=-\tfrac12\,p_R\cdot p_R ,
\end{equation}
and the moduli space of such theories is
$O(\Gamma^{n_L,n_R})\backslash O(n_L,n_R)/(O(n_L)\times O(n_R))$: the position of $\Pi$ fixes
the split of every charge into $p_L$ and $p_R$, and the lattice isometries merely relabel the
states \source{Asp §3.5, eqs.~(77)--(78), ll.~1981--2001}. This is \cref{ch:narain} in
coordinate-free form. Note the convention: with Aspinwall we take the form positive on
left-movers, and we write $|p_R|^2=-p_R\cdot p_R\ge0$. The identity
$p\cdot p=p_L^2-|p_R|^2$ is the level-matching combination $2\,n\cdot m$ of
\cref{ch:narain}.\index{Narain lattice!of signature $(4,20)$}

To recover $G$, $B$ and Wilson lines, choose a null $n_L$-plane $W$ spanned by lattice
vectors, its dual null plane $W^*$, and the remaining negative-definite space
$V\cong\R^{n_R-n_L}$, so $\R^{n_L,n_R}=W\oplus W^*\oplus V$. Then
\begin{equation}\label{eq:stringsk3:graph}
  \Pi=\{\,x+\psi(x)+A(x)\;:\;x\in W\,\},\qquad \psi:W\to W^*,\quad A:W\to V,
\end{equation}
and $\psi$, viewed as a bilinear form on $W$, splits into its symmetric part $G$ and
antisymmetric part $B$. The torus is $W/(\Gamma\cap W)$ with metric $G$; $A$ are the Wilson
lines of the gauge group $U(1)^{n_R-n_L}=V/(\Gamma\cap V)$
\source{Asp §3.5, eqs.~(79)--(84), ll.~2003--2037}. The parameter count is
$\dim\mathrm{Hom}(W,W^*\oplus V)=n_L\cdot n_R$, as it must be by \eqref{eq:stringsk3:dimpq}.

\subsection{Worked example: the circle as a line in \texorpdfstring{$\Gamma^{1,1}$}{Gamma(1,1)}}
\label{sec:stringsk3:circle}
Take $n_L=n_R=1$, $\Gamma^{1,1}=U$ with basis $e,e^*$, $e^2=(e^*)^2=0$, $e\cdot e^*=1$. A
positive line is $\Pi=\R\,(e+t\,e^*)$ with $t>0$, since $(e+te^*)^2=2t$. The unit vector along
$\Pi$ is $u=(e+te^*)/\sqrt{2t}$ and the unit (norm $-1$) vector along $\Pi^\perp$ is
$u_\perp=(e-te^*)/\sqrt{2t}$. For the state $p=n\,e+m\,e^*$:
\begin{equation}\label{eq:stringsk3:circlepLpR}
  p_L^2=(p\cdot u)^2=\frac{(nt+m)^2}{2t}=\tfrac12n^2t+nm+\frac{m^2}{2t},\qquad
  p_R\cdot p_R=-(p\cdot u_\perp)^2=-\frac{(nt-m)^2}{2t},
\end{equation}
which reproduces \source{Asp §3.5, eq.~(88), ll.~2160--2165} with $t=\tan\theta$. Compare
with the book's convention (\cref{ch:circle}): the dimensionless left momentum is
$\sqrt{\ap/2}\,(n_{\rm mom}/R+w_{\rm wind}R/\ap)$, whose square is
$(n_{\rm mom}+w_{\rm wind}\,t)^2/(2t)$ with
\begin{equation}
  t=R^2/\ap .
\end{equation}
So Aspinwall's $m$ is our momentum number, his $n$ is our winding number, and the lattice
isometry $e\leftrightarrow e^*$ maps $t\mapsto1/t$, i.e.\ $R\mapsto\ap/R$: T-duality
(\cref{ch:tduality}). \emph{Numbers}: for $R=2\sqrt{\ap}$, $t=4$, the state $(n,m)=(1,-1)$ has
$p_L^2=9/8$, $|p_R|^2=25/8$ and $p^2=p_L^2-|p_R|^2=-2=2nm$. At the self-dual radius $t=1$ the
same state has $p_L=0$ and $|p_R|^2=2$: a \emph{root} of norm $-2$ has become orthogonal to
$\Pi$. This is the mechanism of \cref{sec:stringsk3:enhanced} in its smallest example
(\cref{fig:stringsk3:circle}).

\begin{figure}[ht]
\centering
\begin{tikzpicture}[>=Stealth,scale=0.78]
  \draw[->] (-0.3,0)--(4.3,0) node[right]{$e$ (winding)};
  \draw[->] (0,-2.2)--(0,3.3) node[above]{$e^*$ (momentum)};
  \foreach \x in {0,1,2,3,4} \foreach \y in {-2,-1,0,1,2,3} \fill[black!45] (\x,\y) circle (1.1pt);
  \draw[very thick,blue!60!black] (0,0)--(3.3,3.3) node[right]{$\Pi$ at $t=1$};
  \draw[thick,blue!60!black,dashed] (0,0)--(4.2,1.4) node[right]{$\Pi$ at $t=\tfrac13$};
  \draw[very thick,red!70!black] (0,0)--(2.1,-2.1) node[right]{$\Pi^\perp$ at $t=1$};
  \draw[->,very thick,red!70!black] (0,0)--(1,-1);
  \node[red!70!black,below left] at (1,-1) {$\alpha=e-e^*$};
\end{tikzpicture}
\caption{The Narain lattice $\Gamma^{1,1}$ of a circle. The modulus $t=R^2/\ap$ is the slope
of the positive line $\Pi$. At $t=1$ the root $\alpha=e-e^*$, $\alpha^2=-2$, lies in $\Pi^\perp$:
it has $p_L=0$, $|p_R|^2=2$, and gives a massless charged vector. Orthogonality is with respect
to the hyperbolic form, for which the axes $e,e^*$ are null.}
\label{fig:stringsk3:circle}
\end{figure}

\subsection{The heterotic string on \texorpdfstring{$T^4$}{T4}}
The heterotic string (\cref{ch:ghosts}) is a superstring on one side and a bosonic string on
the other; the bosonic side has $26-10=16$ extra compact bosons. Following Aspinwall we let the
\emph{left}-movers be the supersymmetric side --- the opposite of the usual convention, forced
by the comparison with K3 \source{Asp §3.5, footnote~11, ll.~2007--2011}. On $T^4$ we then have
$n_L=4$, $n_R=4+16=20$, and the moduli space is \eqref{eq:stringsk3:main} on the nose. In
coordinates \eqref{eq:stringsk3:graph}:
\begin{equation}\label{eq:stringsk3:hetcount}
  \underbrace{\tfrac12\cdot4\cdot5}_{G:\ 10}+\underbrace{\tfrac12\cdot4\cdot3}_{B:\ 6}
  +\underbrace{4\times16}_{\text{Wilson lines }A:\ 64}=80 .
\end{equation}
The same $80$ that was $58+22$ on the K3 side is $10+6+64$ here: the two descriptions
cut the same space into entirely different pieces.\index{heterotic string!on $T^4$}

\section{Type IIA on K3 and the duality with the heterotic string}\label{sec:stringsk3:duality}

\subsection{Supersymmetry in six dimensions}
We now pass from the worldsheet to spacetime: type~IIA string theory on
$\R^{1,5}\times X$. How much supersymmetry survives? Under
$\mathfrak{so}(1,9)\supset\mathfrak{so}(1,5)\oplus\mathfrak{so}(4)$ and
$\mathfrak{so}(4)\cong\mathfrak{su}(2)\oplus\mathfrak{su}(2)$ the ten-dimensional spinor
decomposes as
\begin{equation}\label{eq:stringsk3:spinor}
  \mathbf{16}\to(\mathbf4,\mathbf2,\mathbf1)\oplus(\mathbf4',\mathbf1,\mathbf2).
\end{equation}
The holonomy $SU(2)$ of K3 is one of the two $\mathfrak{su}(2)$ factors, say the second. A
supercharge survives if it is invariant under holonomy (it must be covariantly constant), so
only $(\mathbf4,\mathbf2,\mathbf1)$ survives: $8$ of the $16$ real components, one minimal
spinor in six dimensions. \emph{K3 preserves half of the supersymmetry that a torus would}
\source{Asp §4.1, eq.~(89), ll.~2268--2289}. Type~IIA has two ten-dimensional supercharges of
opposite chirality, $32$ real components; on K3 we keep $16$, and the opposite chiralities
give $N=(1,1)$ in six dimensions \source{Asp §4.2, ll.~2333--2336}. The heterotic string has $N=(1,0)$ in ten dimensions:
$16$ supercharges, all preserved by the flat $T^4$, again $N=(1,1)$ in six dimensions \source{Asp §4.2, ll.~2374--2376}.\index{supersymmetry!in six dimensions}

\subsection{The moduli space of the string theory}
The holonomy of the moduli space of an $N=(1,1)$ theory in six dimensions is
$\mathfrak{sp}(1)\oplus\mathfrak{sp}(1)\cong\mathfrak{so}(4)$. Scalars sit in two kinds of
multiplets: the gravity multiplet contains one real scalar, the dilaton; each matter (vector)
multiplet contains four scalars in the $\mathbf4$ of $\mathfrak{so}(4)$. Hence the
Teichm\"uller space is $O(4,m)/(O(4)\times O(m))\times\R$ with $m$ the number of matter
multiplets \source{Asp §4.2, eq.~(91), ll.~2337--2346}. Since the R--R fields give no scalars
(pitfall box above), comparison with \cref{thm:stringsk3:global} gives $m=20$ and
\begin{equation}\label{eq:stringsk3:MIIA}
  \mathcal M_{\rm IIA}\cong\mathcal M_\sigma\times\R ,\qquad \dim=80+1=81,
\end{equation}
the conformal field theories sitting at the boundary $\Phi\to-\infty$ of weak coupling
\source{Asp §4.2, eq.~(92), ll.~2354--2362}. The form $O(4,m)$ is obtained here from spacetime supersymmetry alone, independently of
\cref{sec:stringsk3:count}; only $m=20$ is borrowed. Now $4$ is the number of scalars in a
multiplet and $20$ the number of multiplets.

The gauge fields come from the R--R sector: the one-form gives one $U(1)$; the three-form
$C_3$ reduced on the $22$ two-cycles gives $22$; and the six-dimensional dual of $C_3$
($\dd C_3^*=*\dd C_3$, a one-form in six dimensions) gives one more:
\begin{equation}\label{eq:stringsk3:24}
  1+22+1=24=b_0+b_2+b_4=\rk H^*(X,\Z)=\rk\Gamma^{4,20}
\end{equation}
\source{Asp §4.3, eq.~(94), ll.~2434--2440}. In the standard multiplet structure $4$ of the $24$ vectors belong to the gravity
multiplet and $20$ to the matter multiplets (standard; not checked against a source in this
book's library). For the heterotic string on $T^4$ the generic
gauge group is also $U(1)^{24}$: $16$ from the extra right-moving bosons, $4$ Kaluza--Klein
vectors from the metric, $4$ from the $B$-field ($b_1(T^4)=4$)
\source{Asp §4.3, ll.~2444--2449}. The charge lattice is $\Gamma^{4,20}$ in both cases: the
Narain lattice on one side, $H^*(X,\Z)$ on the other (\cref{ch:mukai}).

\subsection{The duality statement}
\begin{conjecture}[String--string duality; Aspinwall's Proposition~1, \TL]
\label{conj:stringsk3:duality}
The type~IIA string compactified on a K3 surface is equivalent to the heterotic string
compactified on $T^4$. The moduli spaces $\mathcal M_\sigma\times\R$ are identified in the
obvious way on $\mathcal M_\sigma$, while strong coupling of one theory is mapped to weak
coupling of the other:
\begin{equation}\label{eq:stringsk3:dilaton}
  \Phi_{\rm Het}=-\Phi_{\rm IIA},\qquad g_{6,\rm Het}=\ee^{2\Phi_{\rm Het}}\,g_{6,\rm IIA},
\end{equation}
where $g_6$ is the six-dimensional spacetime metric and $\Phi$ the dilaton
\source{Asp §4.2, Proposition~1 and eq.~(93), ll.~2377--2403}.
\end{conjecture}\index{string--string duality}\index{heterotic/type IIA duality}

We call it a conjecture because our source is explicit that it cannot be proved with the
present definition of string theory; it is supported by the agreement of the moduli spaces,
of the six-dimensional effective actions, and of the soliton spectra, and ``nothing has been
discovered \dots to disprove'' it \source{Asp §4.2, ll.~2391--2413}. It is a statement in the
literature (Tier~\TL\ as a \emph{conjecture with strong evidence}), not a theorem. Notice what
is \emph{not} claimed: the two worldsheet theories are not equivalent conformal field theories.
The $80$-dimensional spaces agree, but the map must reverse the sign of the dilaton, so that
the perturbative expansion of one side is invisible in the other
\source{Asp §4.2, ll.~2377--2388}.

The rescaling in \eqref{eq:stringsk3:dilaton} has a simple consequence. A tension is a mass
per length, i.e.\ it scales as (length)$^{-2}$. If $g_{\rm Het}=\ee^{2\Phi_{\rm Het}}g_{\rm IIA}$,
lengths measured in the heterotic metric are $\ee^{\Phi_{\rm Het}}$ times IIA lengths, so an
object of tension $T$ in IIA units has tension
\begin{equation}\label{eq:stringsk3:tension}
  T^{(\rm Het\ units)}=\ee^{-2\Phi_{\rm Het}}\,T=\frac{T}{\lambda_{\rm Het}^2} .
\end{equation}
The fundamental IIA string, $T=1/(2\pi\ap)$, is therefore a heavy object of tension
$\propto1/\lambda_{\rm Het}^2$ in the weakly coupled heterotic theory --- the characteristic
coupling dependence of a soliton --- and vice versa. Each string is a soliton of the other;
see the references in \source{Asp §4.2, ll.~2404--2409}.

\begin{historybox}
The local form of the moduli space of $N=(4,4)$ theories is due to Seiberg (1988); the global
form with the lattice $\Gamma^{4,20}$ was obtained by Aspinwall and Morrison. The toroidal
moduli space is due to Narain (1986) and to Narain and collaborators (1987). The duality
between type~IIA on K3 and the heterotic string on $T^4$ was suggested by Hull and Townsend
and by Witten (1995), and the appearance of enhanced gauge symmetry at orbifold points of K3
was first shown by Witten in the same paper \source{Asp, reference list [50], [52], [61], [62],
[70], [71], ll.~5355--5398; §3.2, ll.~1671--1673; §3.3, ll.~1731--1732; §4.3, l.~2524}.
\end{historybox}

\section{Enhanced gauge symmetry}\label{sec:stringsk3:enhanced}

\subsection{Heterotic side: roots orthogonal to \texorpdfstring{$\Pi$}{Pi}}
At special points of a Narain moduli space extra massless vectors appear
(\cref{ch:tduality}: $SU(2)\times SU(2)$ at the self-dual radius). In the heterotic string a
spacetime vector is created by a vertex operator whose supersymmetric (left) side carries the
vector index, $\partial X^\mu$ or its superpartner, and whose bosonic (right) side is any
operator of weight $\bar h=1$. Taking the right-moving operator to be a pure exponential of
charge $p$, \eqref{eq:stringsk3:weights} requires $-\tfrac12p_R\cdot p_R=1$, and masslessness
requires that no weight comes from the left momentum, $p_L=0$. Hence
\begin{equation}\label{eq:stringsk3:roots}
  \mathcal A=\big\{\alpha\in\Gamma^{4,20}\cap\Pi^\perp\;:\;\alpha\cdot\alpha=-2\big\}
\end{equation}
is the set of extra massless vectors \source{Asp §4.3, eq.~(95), ll.~2465--2489}. Their charges
under $U(1)^{24}$ are the coordinates of $\alpha$, so $\mathcal A$ is the root system of the
non-abelian part of the gauge group, with Cartan subalgebra inside $\mathfrak u(1)^{24}$. Three
consequences follow at once: the rank stays $24$; all roots have the same length, so the group
is simply laced, of type $A$--$D$--$E$; and since $\Pi^\perp$ is negative-definite of dimension
$20$, the non-abelian part has rank at most $20$.
\index{enhanced gauge symmetry}\index{ADE classification}

Why only norm $-2$ and not $+2$? On the supersymmetric side the GSO projection prevents new
left-moving states from becoming massless \source{Asp §4.3, ll.~2456--2461}. In the bosonic
string on a circle both $\alpha=e-e^*$ (norm $-2$, $p_L=0$) and $\beta=e+e^*$ (norm $+2$,
$p_R=0$) give massless vectors at $t=1$, hence $SU(2)\times SU(2)$; the heterotic string keeps
only the first.

\paragraph{Mass of the $W$-bosons near the special point.} Level matching fixes the mass of
these states away from the enhanced point. We use the standard heterotic mass-shell
conditions, $\tfrac{\ap}4M^2=\tfrac12p_L^2+N_L-\tfrac12=\tfrac12|p_R|^2+N_R-1$, with the
normal-ordering constants $-\tfrac12$ (NS sector) and $-1$ (bosonic) of
\cref{ch:quantum,ch:ghosts} (standard; not checked against a source in this book's library).
For a vector, $N_L=\tfrac12$ and $N_R=0$. Both expressions agree precisely when
$p_L^2-|p_R|^2=-2$, i.e.\ $\alpha^2=-2$, and then
\begin{equation}\label{eq:stringsk3:Wmass}
  M^2=\frac{2}{\ap}\,p_L(\alpha)^2 ,
\end{equation}
where $p_L(\alpha)$ is the orthogonal projection of $\alpha$ to $\Pi$. The mass is the
``distance'' of the root from $\Pi^\perp$: the Higgs mechanism, drawn in a lattice.

\begin{example}[$E_8\times E_8$ without Wilson lines]\label{ex:stringsk3:e8e8}
Split $\Gamma^{4,20}\cong\Gamma^{4,4}\oplus E_8(-1)\oplus E_8(-1)$ and take
$\Pi\subset\Gamma^{4,4}\otimes\R$, i.e.\ $A=0$ in \eqref{eq:stringsk3:graph}. Every vector of
$E_8(-1)\oplus E_8(-1)$ is orthogonal to $\Pi$, in particular all $240+240$ roots of norm $-2$.
The gauge group contains $E_8\times E_8$, of dimension $248+248$, and the generic
$\Gamma^{4,4}$ part adds $U(1)^8$: total rank $8+8+8=24$. The largest possible enhancement has
non-abelian rank $20$, for instance $E_8\times E_8\times SU(2)^4\times U(1)^4$
\source{Asp §4.3, ll.~2490--2499 and 2551--2554}.
\end{example}

\subsection{Type IIA side: shrinking two-spheres}
In type~IIA perturbation theory the gauge fields are R--R fields, under which no perturbative
string state is charged; the conformal field theory can only ever see $U(1)^{24}$
\source{Asp §4.3, ll.~2440--2443}. If \cref{conj:stringsk3:duality} holds, the non-abelian gauge
symmetry must nevertheless be there. Where?

Choose $w$ as in \cref{sec:stringsk3:fourplane}; from now on $\alpha$ denotes a root, and the
coefficient called $\alpha$ in \eqref{eq:stringsk3:Bprime} is written $\alpha_{\rm vol}$. Suppose
all roots $\alpha\in\mathcal A$ lie
in $w^\perp/w=H^2(X,\Z)$. Use the basis \eqref{eq:stringsk3:basis} of $\Pi$. Since
$\alpha\cdot w=\alpha\cdot w^*=0$,
\begin{equation}\label{eq:stringsk3:conditions}
  \alpha\cdot\xi_i=\alpha\cdot\sigma_i,\qquad \alpha\cdot B'=\alpha\cdot B ,
\end{equation}
so $\alpha\perp\Pi$ is equivalent to \emph{two} conditions:
\begin{enumerate}[leftmargin=*]
\item $\alpha\perp\Sigma$. A class with $\alpha^2=-2$ orthogonal to the holomorphic two-form
 lies in the Picard lattice and $\pm\alpha$ is the class of a rational curve $C\cong S^2$; if
 moreover $J\cdot\alpha=0$ this curve has zero area, and $X$ has an orbifold singularity there
 \source{Asp §2.6, Theorem~4, ll.~1196--1211} (\cref{ch:kummer}). The roots orthogonal to
 $\Sigma$ form the root system of the $A$--$D$--$E$ singularity: $\C^2/\Z_n$ corresponds to
 $A_{n-1}$, i.e.\ $SU(n)$.
\item $\alpha\cdot B=0$: the $B$-field has no component along the vanishing cycle.
\end{enumerate}

\begin{proposition}[Aspinwall's Proposition~2 (Witten), \TL, conditional on
\cref{conj:stringsk3:duality}]\label{thm:stringsk3:ADE}
A type~IIA string on a K3 surface with an orbifold singularity of type $A$--$D$--$E$, with
vanishing $B$-field through the collapsed cycles, has non-abelian gauge symmetry with the
simply-laced Lie algebra of the same $A$--$D$--$E$ type
\source{Asp §4.3, Proposition~2, ll.~2523--2534}.
\end{proposition}

The volume $V$ does not enter \eqref{eq:stringsk3:conditions}: the enhancement happens for a
K3 of any size. The physical mechanism, outside perturbation theory, is that R--R charged
solitons --- D2-branes in the language of \cref{ch:dbranes} --- wrap the two-sphere; their mass
is proportional to its area and they become massless when it shrinks
\source{Asp §4.3, ll.~2555--2563}. They are the $W$-bosons.
\index{D-brane!wrapped on a vanishing cycle}

\begin{figure}[ht]
\centering
\begin{tikzpicture}[>=Stealth,scale=0.95]
  % IIA side
  \node[draw,rounded corners,align=center,inner sep=5pt] (A) at (0,0)
     {type IIA on K3\\[2pt]\small $58$ metric $+\,22$ $B$-field\\\small gauge fields: R--R, $1+22+1$};
  \node[draw,rounded corners,align=center,inner sep=5pt] (H) at (8.4,0)
     {heterotic on $T^4$\\[2pt]\small $10\ (G)+6\ (B)+64$ (Wilson)\\\small gauge fields: $16+4+4$};
  \node[draw,rounded corners,align=center,inner sep=5pt,fill=yellow!10] (M) at (4.2,-2.3)
     {$O(\Gamma^{4,20})\backslash O(4,20)/(O(4)\times O(20))\;\times\;\R_{\Phi}$\\[2pt]
      \small charge lattice $\Gamma^{4,20}$, generic gauge group $U(1)^{24}$};
  \draw[<->,thick] (A)--(H) node[midway,above]{\small $\Phi_{\rm Het}=-\Phi_{\rm IIA}$};
  \draw[->] (A)--(M); \draw[->] (H)--(M);
  \node[align=center,font=\small] at (0,-3.8) {$S^2$ shrinks, $B\!\cdot\!\alpha=0$:\\wrapped brane massless};
  \node[align=center,font=\small] at (8.4,-3.8) {root $\alpha\perp\Pi$, $\alpha^2=-2$:\\winding state massless};
  \draw[<->,dashed] (2.3,-3.8)--(6.0,-3.8) node[midway,above]{\small same point of $\mathcal M_\sigma$};
\end{tikzpicture}
\caption{The two descriptions of the six-dimensional $N=(1,1)$ theory. The $80$ moduli are
counted differently on the two sides; the enhanced gauge symmetry is perturbative on the
heterotic side and non-perturbative on the type~IIA side.}
\label{fig:stringsk3:duality}
\end{figure}

\subsection{Worked example: an \texorpdfstring{$A_1$}{A1} singularity and the role of
\texorpdfstring{$B$}{B}}\label{sec:stringsk3:A1}
Let $\alpha\in H^2(X,\Z)$ be the class of a rational curve, $\alpha^2=-2$. The basis
\eqref{eq:stringsk3:basis} is orthogonal with $\xi_i^2=1$ and $\xi_4^2=V$, so by
\eqref{eq:stringsk3:conditions} the projection of $\alpha$ to $\Pi$ has length
\begin{equation}\label{eq:stringsk3:pLalpha}
  p_L(\alpha)^2=\sum_{i=1}^3(\alpha\cdot\sigma_i)^2+\frac{(\alpha\cdot B)^2}{V},
  \qquad V=2\alpha_{\rm vol}+B\cdot B .
\end{equation}
The first term is the squared
``size'' of the curve measured by the three normalized hyperk\"ahler forms; the second is the
$B$-field flux through it.

\emph{(a) Small smooth sphere.} Take $(\alpha\cdot\sigma_i)=(0.06,\,0,\,0.08)$ and
$\alpha\cdot B=0$. Then $p_L^2=0.0036+0.0064=0.01$ and by \eqref{eq:stringsk3:Wmass}
$M^2=0.02/\ap$, $M\approx0.14/\sqrt{\ap}$: a light charged vector pair $\pm\alpha$. At
$\alpha\cdot\sigma_i=0$ the pair is massless and, with the $U(1)$ under which it is charged,
fills out $SU(2)$; the gauge group is $SU(2)\times U(1)^{23}$, of dimension $3+23=26$.

\emph{(b) The orbifold conformal field theory.} An orbifold such as $T^4/\Z_2$
(\cref{ch:kummer}) is a perfectly regular conformal field theory, although its $16$ two-spheres
have zero size. There is no contradiction: the orbifold conformal field theory has
$B$-field flux $\tfrac12$ through each collapsed sphere
\source{Asp §4.3, ll.~2535--2547}, that is $\alpha\cdot B=\tfrac12$ in our notation. Then
\eqref{eq:stringsk3:pLalpha} gives $p_L^2=1/(4V)$ and
\begin{equation}
  M^2=\frac{1}{2V\ap}\neq0 .
\end{equation}
For $V=100$ (a large K3 in units of $\ap{}^2$) this is $M\approx0.07/\sqrt{\ap}$: light, but not
massless. The singular conformal field theory sits at $\alpha\cdot B=0$, half a period away.
Because $B\sim B+e$, the flux $\alpha\cdot B$ is periodic with period one, and $\tfrac12$ is the
point farthest from the enhancement locus.
\index{orbifold!B-field at the orbifold point@$B$-field at the orbifold point}

\emph{(c) A stringy enhancement at small volume.} Roots need not lie in $H^2$. Take
$\alpha=w^*-w$: $\alpha^2=-2\,w\cdot w^*=-2$. Now $\alpha\cdot\xi_i=-(w^*-w)\cdot w\,
(\sigma_i\cdot B)=-\sigma_i\cdot B$ and $\alpha\cdot B'=\alpha_{\rm vol}-1$. So $\alpha\perp\Pi$
requires $B\perp\Sigma$ and $\alpha_{\rm vol}=1$; with $B=0$ this is $V=2$. A \emph{smooth} K3
surface of volume of order one in string units has enhanced $SU(2)$ symmetry, whatever its
shape. This is the K3 counterpart of the self-dual radius of \cref{sec:stringsk3:circle}, where
the root $e-e^*$ also pairs winding with momentum; Aspinwall states the general conclusion,
that if the roots are not all orthogonal to $w$ the surface may be smooth but its volume is
of order one \source{Asp §4.3, ll.~2541--2550}. (The specific value $V=2$ depends on the
normalization \eqref{eq:stringsk3:vol}; the computation here is ours.)

\begin{pitfallbox}[Common pitfall: ``singular K3 $\Rightarrow$ non-abelian gauge symmetry'']
An orbifold singularity is necessary (for roots in $H^2$) but not sufficient: the $B$-field
through the vanishing cycle must also vanish. Conversely a smooth, small K3 can have enhanced
symmetry, example (c). The locus of enhanced symmetry has real codimension \emph{four} in
$\mathcal M_\sigma$ --- three conditions on the metric and one on $B$, matching the four scalars
of the vector multiplet that becomes part of the non-abelian theory.
\end{pitfallbox}

\begin{remark}[Type IIB on K3]\label{rem:stringsk3:IIB}
For type~IIB the six-dimensional theory is chiral, $N=(2,0)$, the holonomy is
$\mathfrak{sp}(2)\cong\mathfrak{so}(5)$, and the R--R potentials of even degree \emph{do}
give scalars: $b_0+b_2+b_4=24$. With the dilaton the count is $58+22+1+1+22+1=105=5\times21$
and the moduli space is $O(\Gamma^{5,21})\backslash O(5,21)/(O(5)\times O(21))$, in which the
dilaton is no longer a separate factor \source{Asp §4.4, eq.~(97) and table, ll.~2573--2621}.
The pattern $(3,19)\to(4,20)\to(5,21)$ continues on $\KT$ with the lattice of signature
$(6,22)$ (\cref{ch:k3t2}).
\end{remark}

\section{What the machine has checked}\label{sec:stringsk3:lean}

No part of this chapter's physics --- sigma models, moduli spaces as quotients, string
duality --- is formalized. What is formalized are the lattice-theoretic facts on which the
argument rests, in the library \lean{DualScaleStream2.Lattice}, and a handful of numerical
identities in the older library \lean{StringTheoryFoundation}. We list them with their exact
scope.

\subsection{Signatures and ranks}
The file \lean{DualScaleStream2/Lattice/K3T2Signature.lean} defines a signature as a pair of
natural numbers with componentwise addition, takes $\sgn U=(1,1)$ and
$\sgn E_8(-1)=(0,8)$ as \emph{inputs}, and computes.

\begin{leanbox}[In Lean: $(3,19)+(1,1)=(4,20)$ and the mod-8 condition]
\begin{leancode}
def sigK3 : Signature := sigE8Neg + sigE8Neg + sigU + sigU + sigU
def sigMukai : Signature := sigK3 + sigU

theorem sigMukai_eq : sigMukai = ⟨4, 20⟩ := by ...

theorem index_mod_eight :
    sigK3.index % 8 = 0 ∧ sigMukai.index % 8 = 0 ∧
    sigK3T2.index % 8 = 0 := by ...
\end{leancode}
\textbf{Reading.} Adding the pairs $(0,8)+(0,8)+3\cdot(1,1)+(1,1)$ gives $(4,20)$, and
$4-20=-16$ is divisible by $8$ (as are $3-19$ and $6-22$). \textbf{Proof idea.} Both sides
reduce to the same numerals by unfolding definitions (\lean{rfl}). \textbf{Scope.} This is
bookkeeping on pairs of numbers. That $U$ and $E_8(-1)$ \emph{have} these signatures, that
$H^*(K3,\Z)$ \emph{is} $4U\oplus2E_8(-1)$, and the theorem that even unimodular lattices exist
only for $n_+-n_-\equiv0\pmod 8$ --- which is what forces $n_L-n_R\in8\Z$ in
\cref{sec:stringsk3:narain} --- are \TL\ inputs. The companion theorems \lean{sigK3\_eq},
\lean{rank\_K3} ($22$) and \lean{sigK3\_matches\_hodge} were discussed in \cref{ch:k3lattice}.
\end{leanbox}

\subsection{Evenness of the Mukai lattice and the \texorpdfstring{$B$}{B}-shift}
In \cref{thm:stringsk3:bshift} integrality of $T_e$ rested on $e\cdot e\in2\Z$. The statement
that adjoining $H^0\oplus H^4$ preserves evenness is kernel-checked for an arbitrary even Gram
matrix (\cref{ch:mukai}):

\begin{leanbox}[In Lean: the extended lattice is even if $H^2$ is]
\begin{leancode}
def mukaiPair {n : ℕ} (L : Gram n) (v w : MukaiVec n) : ℤ :=
  v.c ⬝ᵥ (L *ᵥ w.c) - v.r * w.s - v.s * w.r

theorem mukaiPair_even {n : ℕ} (L : Gram n) (hL : Lᵀ = L) (heven : IsEvenDiag L)
    (v : MukaiVec n) : Even (mukaiPair L v v) := by ...

theorem hyperbolicUNeg_unimodular : IsUnimodular hyperbolicUNeg := ...
\end{leancode}
\textbf{Reading.} For a symmetric integer matrix $L$ with even diagonal and any integer
vector $v=(r,c,s)$, the number $c^{\mathsf T}Lc-2rs$ is even; and the $2\times2$ Gram matrix
$\left(\begin{smallmatrix}0&-1\\-1&0\end{smallmatrix}\right)$ of $H^0\oplus H^4$ has determinant
$\pm1$. \textbf{Proof idea.} $c^{\mathsf T}Lc=\sum_iL_{ii}c_i^2+2\sum_{i<j}L_{ij}c_ic_j$ is even
by the general lemma \lean{even\_quadratic\_form\_of\_even\_diag}, and $2rs$ is even.
\textbf{Scope.} Note the sign: the Mukai pairing uses $U(-1)$ on $H^0\oplus H^4$, Aspinwall's
$w,w^*$ span $U$. Since $U(-1)\cong U$ (send $w^*\mapsto-w^*$) this is a convention, not a
discrepancy, but the isomorphism is not a Lean theorem. The map $T_e$ of
\eqref{eq:stringsk3:Te} itself is \emph{not} formalized; see \cref{exo:stringsk3:leanTe}.
\end{leanbox}

\subsection{Roots and reflections}
The enhanced gauge symmetry of \cref{sec:stringsk3:enhanced} is governed by vectors of norm
$-2$. For each such vector the reflection $s_\alpha(x)=x+(x\cdot\alpha)\alpha$ is an isometry of
the lattice; these reflections generate the Weyl group of the enhanced gauge algebra, and on
the K3 side they are the Picard--Lefschetz monodromies around the singular surface
(\cref{ch:k3lattice}). The algebraic core is Tier~\TA:

\begin{leanbox}[In Lean: reflections in $(-2)$-vectors are integral isometries]
\begin{leancode}
def reflection (L : Gram n) (v : Fin n → ℤ) : Gram n :=
  1 + vecMulVec v v * L

theorem reflection_isometry (L : Gram n) (hL : Lᵀ = L) (v : Fin n → ℤ)
    (hv : latticeNorm L v = -2) :
    (reflection L v)ᵀ * L * reflection L v = L := by ...

theorem e8Neg_simpleRoot_norm (i : Fin 8) :
    latticeNorm e8Neg (Pi.single i 1) = -2 := by ...
\end{leancode}
\textbf{Reading.} For any symmetric integer Gram matrix $L$ and any integer vector $v$ with
$v^{\mathsf T}Lv=-2$, the integer matrix $S=1+vv^{\mathsf T}L$ satisfies $S^{\mathsf T}LS=L$.
The eight basis vectors of $E_8(-1)$ have norm $-2$. \textbf{Proof idea.} Expand
$S^{\mathsf T}LS=L+2Lvv^{\mathsf T}L+(v^{\mathsf T}Lv)\,Lvv^{\mathsf T}L$ and insert $-2$.
\textbf{Scope.} This is the statement used in \cref{ex:stringsk3:e8e8} only for the eight simple
roots of one $E_8(-1)$; the count of $240$ roots, and the identification of $\mathcal A$ with a
gauge algebra, are \TL. (Declarations in
\lean{DualScaleStream2/Lattice/Reflection.lean}.)
\end{leanbox}

\subsection{Arithmetic shadows of the duality}
The older file \lean{StringTheoryFoundation/StringTheory/WittenDuality.lean} is named after the
duality but contains no string theory. Its four theorems are:

\begin{leanbox}[In Lean: four numerical identities --- read the statements, not the names]
\begin{leancode}
def wittenModuliDimension (p q : Nat) : Nat := p * q
theorem witten_moduli_dim_is_80 :
    wittenModuliDimension 4 20 = 80 := by decide

theorem witten_duality_rank_match :
    defaultDualityLattice.b2_K3 + 2
      = defaultDualityLattice.t4_momentum + defaultDualityLattice.t4_gauge := by
  decide

theorem witten_6d_supercharges :
    defaultSixDSusy.numSupercharges = 16 := by decide

def exhibitsGaugeEnhancement (c : TwoCycle) : Prop :=
  c.intersectionSelf = -2 ∧ c.area = 0
theorem witten_gauge_enhancement_at_singularity :
    exhibitsGaugeEnhancement ⟨-2, 0⟩ := by ...
\end{leancode}
\textbf{Reading.} (i) $4\cdot20=80$. (ii) For a record with default fields $22$, $4$, $20$:
$22+2=4+20$. (iii) A record field defined as $32/2$ equals $16$. (iv) The pair $(-2,0)$
satisfies the predicate ``first component is $-2$ and second component is $0$''.
\textbf{Scope.} These are arithmetic shadows of \eqref{eq:stringsk3:dimpq},
\eqref{eq:stringsk3:24} and of \cref{sec:stringsk3:enhanced}. In (iv) the word ``gauge
enhancement'' is the \emph{name} of a predicate on two numbers; nothing about gauge fields is
proved, and the $B$-field condition of \cref{thm:stringsk3:ADE} is absent from the predicate.
The informal comments in that file also describe the $22$ extra moduli as Ramond--Ramond
periods and include the dilaton in the $80$, which contradicts our source (pitfall box in
\cref{sec:stringsk3:count}); comments are not checked by the kernel. A faithful formalization
would need at least the Grassmannian of positive four-planes in a real quadratic space of
signature $(4,20)$ with $O(4U\oplus2E_8(-1))$ acting on it;
\cref{thm:stringsk3:plane,thm:stringsk3:bshift} are elementary linear algebra and realistic
first targets (\cref{ch:open}).
\end{leanbox}

\begin{table}[ht]
\centering\small
\begin{tabular}{@{}p{0.45\textwidth}p{0.34\textwidth}l@{}}
\toprule
Claim & Where & Tier\\
\midrule
$(3,19)+(1,1)=(4,20)$; index $\equiv0 \pmod 8$ at $(3,19),(4,20),(6,22)$ &
  \lean{sigMukai\_eq}, \lean{index\_mod\_eight} & \TA\\
$H^2$ even $\Rightarrow$ $H^0\oplus H^2\oplus H^4$ even; $U(-1)$ unimodular &
  \lean{mukaiPair\_even}, \lean{hyperbolicUNeg\_unimodular} & \TA\\
Reflection in a $(-2)$-vector is an integral isometry; simple roots of $E_8(-1)$ have norm $-2$ &
  \lean{reflection\_isometry}, \lean{e8Neg\_simpleRoot\_norm} & \TA\\
$4\cdot20=80$; $22+2=4+20$; $32/2=16$ &
  \lean{witten\_moduli\_dim\_is\_80}, \lean{witten\_duality\_rank\_match},
  \lean{witten\_6d\_supercharges} & \TA\ (shadow)\\
\Cref{thm:stringsk3:plane,thm:stringsk3:bshift}; formulas \eqref{eq:stringsk3:Wmass},
  \eqref{eq:stringsk3:pLalpha}; example (c) & derived here; sympy check & ---\\
K3 sigma model is $N=(4,4)$, metric exact & Asp §3.1 & \TL\\
$\mathcal M_\sigma=O(\Gamma^{4,20})\backslash O(4,20)/(O(4)\times O(20))$ (under stated
  assumptions) & Asp §3.2--3.4, Thms.~5, 6 & \TL\\
$V=2\alpha+B^2$ & Asp §3.3, eq.~(68) & \TL\\
IIA on K3 $\leftrightarrow$ heterotic on $T^4$, $\Phi\mapsto-\Phi$ & Asp §4.2, Prop.~1 &
  \TL\ (conjecture)\\
$A$--$D$--$E$ enhanced gauge symmetry at orbifold points with $B\cdot\alpha=0$ & Asp §4.3,
  Prop.~2 & \TL\ (conditional)\\
Identification of any Lean object above with a physical system & --- & never \TA\\
\bottomrule
\end{tabular}
\caption{Tier table for this chapter. The \TA\ entries depend only on the axioms
\lean{propext}, \lean{Classical.choice}, \lean{Quot.sound}. Rows marked ``derived here'' are
ordinary pen-and-paper mathematics of this chapter, checked symbolically but not
formalized.}
\label{tab:stringsk3:tiers}
\end{table}

\begin{summarybox}
\begin{itemize}[leftmargin=*]
\item A string on a K3 surface is described by an $N=(4,4)$ superconformal sigma model; the
  Ricci-flat metric is an exact solution. Its parameters are the metric ($58$) and the class
  of the $B$-field ($22$): $80=4\times20$ in all. The dilaton is not among them.
\item Locally the moduli space is the Grassmannian $O(4,20)/(O(4)\times O(20))$ of positive
  four-planes $\Pi\subset\R^{4,20}$; globally one divides by $O(\Gamma^{4,20})$, the isometry
  group of the Mukai lattice $H^*(X,\Z)$ (under completeness and Hausdorff assumptions).
\item After choosing a null vector $w$, $\Pi$ is spanned by $\sigma_i-(\sigma_i\cdot B)w$ and
  $B'=\alpha w+w^*+B$, with $B'\cdot B'=2\alpha+B^2=V$. Integral $B$-shifts are lattice isometries
  because the lattice is even. Changing $w$ is a T-duality-like relabelling.
\item The same coset is the Narain moduli space of the heterotic string on $T^4$
  ($10+6+64=80$); both theories have $N=(1,1)$ supersymmetry in six dimensions and generic
  gauge group $U(1)^{24}$. This supports the conjectural duality IIA/K3 $=$ Het/$T^4$ with
  $\Phi\mapsto-\Phi$.
\item Extra massless vectors are roots $\alpha\in\Gamma^{4,20}\cap\Pi^\perp$, $\alpha^2=-2$;
  the gauge group is simply laced of rank $24$. On the K3 side this means an $A$--$D$--$E$
  orbifold singularity \emph{and} vanishing $B$-flux through the collapsed spheres; the
  orbifold conformal field theory has flux $\tfrac12$ and is regular.
\item Machine-checked: signature and parity bookkeeping, reflections in $(-2)$-vectors,
  numerical identities. Not machine-checked: anything about sigma models, moduli spaces, duality.
\end{itemize}
\end{summarybox}

\section*{Exercises}
\addcontentsline{toc}{section}{Exercises}

\begin{exercise}[$\star$]\label{exo:stringsk3:dims}
Verify \eqref{eq:stringsk3:dimpq} and evaluate it for $(3,19)$, $(4,20)$, $(5,21)$ and $(6,22)$.
Match the first three numbers with a count of physical fields
(\eqref{eq:stringsk3:einstein} without the volume, \eqref{eq:stringsk3:80},
\cref{rem:stringsk3:IIB}).
\end{exercise}

\begin{exercise}[$\star$]\label{exo:stringsk3:circle}
For the circle of \cref{sec:stringsk3:circle} with $t=R^2/\ap=9$, compute $p_L^2$ and $|p_R|^2$
for $(n,m)=(1,-1)$, $(1,1)$ and $(0,1)$, and check $p_L^2-|p_R|^2=2nm$ in each case. Which of the
three states is a root? Find all $t$ at which some root of $\Gamma^{1,1}$ is orthogonal to $\Pi$.
\end{exercise}

\begin{exercise}[$\star$]\label{exo:stringsk3:gauge}
List the $24$ abelian gauge fields of type~IIA on K3 and of the heterotic string on $T^4$ by
their ten-dimensional origin. Which Lean theorem of \cref{sec:stringsk3:lean} records the
equality of the two counts, and what exactly does it state?
\end{exercise}

\begin{exercise}[$\star\star$]\label{exo:stringsk3:group}
Show that the maps $T_e$ of \cref{thm:stringsk3:bshift} satisfy $T_e\circ T_f=T_{e+f}$, and that
for $g\in O(\Gamma^{3,19})$ (extended by the identity on $w,w^*$) one has
$g\,T_e\,g^{-1}=T_{g(e)}$. Conclude that the $T_e$ and $O(\Gamma^{3,19})$ generate the semidirect
product $O(\Gamma^{3,19})\ltimes\Gamma^{3,19}$.
\end{exercise}

\begin{exercise}[$\star\star$]\label{exo:stringsk3:A2}
Let $\alpha_1,\alpha_2\in H^2(X,\Z)$ with $\alpha_1^2=\alpha_2^2=-2$, $\alpha_1\cdot\alpha_2=1$, both
orthogonal to $\Pi$. Show that $\mathcal A$ contains exactly the six vectors
$\pm\alpha_1,\pm\alpha_2,\pm(\alpha_1+\alpha_2)$ among integer combinations of $\alpha_1,\alpha_2$,
identify the gauge group, give its dimension, and name the orbifold singularity of $X$
\source{Asp §2.6, ll.~1227--1240}.
\end{exercise}

\begin{exercise}[$\star\star$, Lean]\label{exo:stringsk3:leansig}
In the style of \lean{sigMukai\_eq}, define \lean{sigIIB : Signature := sigMukai + sigU} and
prove \lean{sigIIB = ⟨5, 21⟩}, \lean{sigIIB.rank = 26} and \lean{sigIIB.index \% 8 = 0}. State in
one sentence what these theorems do and do not say about type~IIB string theory.
\end{exercise}

\begin{exercise}[$\star\star$, Lean]\label{exo:stringsk3:leanroot}
Using \lean{mukaiPair} with an arbitrary \lean{L : Gram n}, prove that the vector
$(r,c,s)=(1,0,1)$ has self-pairing $-2$ (this is \lean{structureSheaf\_mukai\_sq}) and that
$(1,0,-1)$ has self-pairing $+2$. Relate the two vectors to the roots $w^*\mp w$ of example (c)
in \cref{sec:stringsk3:A1}, taking the sign convention $U(-1)$ versus $U$ into account.
\end{exercise}

\begin{exercise}[$\star\star\star$, Lean]\label{exo:stringsk3:leanTe}
Formalize \cref{thm:stringsk3:bshift} for a general symmetric even Gram matrix
\lean{L : Gram n}: write the $(n+2)\times(n+2)$ block Gram matrix of $U\oplus L$, the integer
matrix of $T_e$ (you will need \lean{(e ⬝ᵥ (L *ᵥ e)) / 2} and the hypothesis
\lean{IsEvenDiag L}), and prove $T_e^{\mathsf T}\,(U\oplus L)\,T_e=U\oplus L$. Which lemma of
\cref{sec:stringsk3:lean} supplies the divisibility by two?
\end{exercise}

\begin{exercise}[$\star\star\star$]\label{exo:stringsk3:het}
Work out \eqref{eq:stringsk3:graph} for the heterotic string on a circle with one Wilson line
component switched on: $\Gamma^{1,17}=U\oplus E_8(-1)\oplus E_8(-1)$, $\Pi=\R\,(e+t\,e^*+a)$
with $a\in E_8(-1)\otimes\R$ in the first factor. Show that $\Pi$ is positive iff
$2t+a\cdot a>0$, and that a root $\beta$ of $E_8(-1)$ remains orthogonal to $\Pi$ iff
$\beta\cdot a=0$. For $a$ proportional to a fundamental coweight, which subgroup of $E_8$
survives? (Compare $2t+a\cdot a$ with \eqref{eq:stringsk3:volume}: the Wilson line plays the role
of the $B$-field.)
\end{exercise}

\section*{Further reading}
\addcontentsline{toc}{section}{Further reading}
\begin{itemize}[leftmargin=*]
\item The sigma model, the count $58+22$, and Theorem~5: \source{Asp §3.1--3.2, ll.~1559--1725}.
  The null vector $w$, the decomposition into metric, $B$-field and volume, and the space
  group: \source{Asp §3.3, ll.~1727--1832}. Mirror symmetry for K3 and Theorem~6:
  \source{Asp §3.4, ll.~1836--1963}.
\item Narain moduli spaces as Grassmannians, including the heterotic case and the circle:
  \source{Asp §3.5, ll.~1967--2170}; compare \cref{ch:narain,ch:odd}.
\item Supersymmetry counting, the IIA moduli space and the duality conjecture:
  \source{Asp §4.1--4.2, ll.~2189--2413}. Enhanced gauge symmetry: \source{Asp §4.3,
  ll.~2428--2566}. Type~IIB and $O(5,21)$: \source{Asp §4.4, ll.~2571--2640}.
\item $A$--$D$--$E$ singularities and $(-2)$-curves: \source{Asp §2.6, ll.~1172--1224};
  \cref{ch:kummer}. The Mukai lattice as $H^*(X,\Z)$ and its isometries from derived equivalences:
  \cref{ch:mukai}. What becomes of $\Gamma^{4,20}$ on $\KT$: \cref{ch:k3t2}. Symmetries of K3
  sigma models acting on $\Gamma^{4,20}$ and their relation to $\Mtf$: \cref{ch:moonshine}.
\end{itemize}
